\begin{abstract}
\noindent The saturating-removal (SR) model describes aging as damage that is produced at a rising rate, removed by a
saturating process and perturbed by noise, with death when damage crosses a threshold. It explains the shape of
survival curves across species, and is increasingly used to ask which damage process an intervention or a genetic
difference changes. Answering that question from survival data needs a fitting method whose errors are known. We
describe the methods now used to fit the SR model to survival data, and the evidence for them.

The likelihood is exact: it is obtained by solving the model's Fokker--Planck equation for the first-passage
distribution, and it handles deaths recorded at a finite resolution, censoring, left truncation, a fit window, an
external hazard and parameter changes at a given age. A maximum-likelihood calibrator reports each parameter with a
profile-likelihood interval, as a value, a bound or no claim. Survival data determine the damage production slope,
the noise and the mean removal, but not the saturation constant of removal; we give the reporting rules that follow.

The likelihood agrees with 60-digit arithmetic to \C{acc-mp-v1-lowmass} nats on its hardest rows. In a
benchmark of \C{gate-v6-n} synthetic cohorts with known parameters, nominal 95\,\% intervals for the production slope
covered the truth in \C{v6-pure-eta-J-hits} of \C{v6-pure-eta-J-n} cohorts without external deaths or delayed entry;
with external deaths or a delayed entry, the variant that models them restored coverage. For comparisons between groups we give a naming rule and its measured
error rates: close to nominal for independent individuals, and too permissive when individuals housed together (in cages, vials or plates) are handled by
design-effect scaling instead of a simulated threshold. The code, the frozen calibrator versions, the tests and the
benchmark records are released with pinned versions, so that studies can cite the exact method they used and readers
can recheck the numbers reported here.
\end{abstract}
